\newpage
%%%%%%%%%%%%%%%%%%%%%%%%%%%%%%%%%%%%%%%%%%%%%%%%%%%%%%%%%%%%%%%%%%%%%%%%%%%%%%%%
\section{Memory arenas} %%%%%%%%%%%%%%%%%%%%%%%%%%%%%%%%%%%%%%%%%%%%%%%%%%%%%%%%
%%%%%%%%%%%%%%%%%%%%%%%%%%%%%%%%%%%%%%%%%%%%%%%%%%%%%%%%%%%%%%%%%%%%%%%%%%%%%%%%
\label{sec:arenas}

Many user applications require a different memory allocator than the basic C
\verb'malloc' and \verb'free', and some applications require multiple allocators
to be used at the same time.

GraphBLAS 10.4.0 introduces a new feature to the library to support the use of
multiple allocators: {\em memory arenas}.  An arena is simply a set memory
allocator routines with the same API as the C \verb'malloc' and \verb'free', plus
optional \verb'realloc' and \verb'calloc' methods.

Each object in GraphBLAS lives in one or two memory arenas, depending on the
object.  The \verb'GrB_Matrix', \verb'GrB_Vector', and \verb'GrB_Scalar' can
use two arenas: a {\em header arena} and a {\em data arena}.  Each
object in GraphBLAS is a pointer to a struct with opaque content, typically at
most a few hundred bytes in size, regardless of the size of the matrix.  The
\verb'GrB_Matrix A' is simply a pointer to this struct, which is the header of
the object \verb'A'.  Changing the header of a matrix, vector, or scalar will
change the pointer to the object itself.

The entries in a matrix, vector, or scalar are held in a set of integer arrays
for row/column indices and offsets, plus a numerical array for the data type.
These arrays are all in the data arena of the matrix, vector, or scalar.

All other objects have only a single arena, the header arena.

By default, only a single memory arena is used, and this arena defaults to the
standard C \verb'malloc'/\verb'calloc'/\verb'realloc'/\verb'free' methods.  All
arena features are backward compatible with prior versions of GraphBLAS, so if
no additional arenas are needed, nothing needs to change in applications that
used GraphBLAS versions prior to 10.4.0 (including the MATLAB/Octave
interface).

Objects in different arenas can be mixed arbitrarily, so for any given
GraphBLAS method, not all of its parameters need to refer to objects that are
all in the same arena.

%-------------------------------------------------------------------------------
\subsection{Creating arenas}
%-------------------------------------------------------------------------------

The default arena is initialized when \verb'GrB_init' or \verb'GxB_init' is
called.  \verb'GxB_init' takes in the four pointers to the \verb'malloc',
\verb'calloc', \verb'realloc', and \verb'free' functions.  This becomes arena
0, the default arena (\verb'GrB_DEFAULT', which is defined as zero).

Once an arena is initialized, it cannot be deleted except by a call to
\verb'GrB_finalize', which resets all arenas (including the default arena).
The macro \verb'GxB_NARENAS' is set in \verb'GraphBLAS.h' and defines the
maximum number of arenas (8 in GraphBLAS version 10.4.0).

New arenas can be initialized after \verb'GrB_init' with the
\verb'GxB_arena_init' method.

\begin{mdframed}[userdefinedwidth=6in]
{\footnotesize
\begin{verbatim}
GrB_Info GxB_arena_init     // create a new arena
(
    // input
    int arena,              // 2 to GxB_NARENAS-1
    void * (* user_malloc_function  ) (size_t),         // required
    void * (* user_calloc_function  ) (size_t, size_t), // not used
    void * (* user_realloc_function ) (void *, size_t), // optional, can be NULL
    void   (* user_free_function    ) (void *)          // required
) ;
\end{verbatim}
}\end{mdframed}

Arena 1 is reserved for future use, and arena 0 is only initialized by
\verb'GrB_init' or \verb'GxB_init', so the \verb'arena' parameter in
\verb'GxB_arena_init' can be in the range 2 to \verb'GxB_NARENAS-1'.

Just as in \verb'GxB_init', the \verb'calloc' and \verb'realloc' methods are
optional.  GraphBLAS currently makes no use of the \verb'calloc' method itself,
but it can be returned by \verb'GrB_get'.  It uses the \verb'realloc' method if
avaliable; otherwise if this function is NULL, it uses \verb'malloc' and
\verb'memcpy' in its place.

To retrieve these function pointers after the arena is initialized, use the following:

    {\footnotesize
    \begin{verbatim}
    GrB_Global_get_VOID (GrB_GLOBAL, &malloc_func,  GxB_ARENA_MALLOC  + arena) ;
    GrB_Global_get_VOID (GrB_GLOBAL, &calloc_func,  GxB_ARENA_CALLOC  + arena) ;
    GrB_Global_get_VOID (GrB_GLOBAL, &realloc_func, GxB_ARENA_REALLOC + arena) ;
    GrB_Global_get_VOID (GrB_GLOBAL, &free_func,    GxB_ARENA_FREE    + arena) ; \end{verbatim}}

%-------------------------------------------------------------------------------
\subsection{Using arenas via the global or per-thread context}
%-------------------------------------------------------------------------------

Once an arena is initialized, it can be used to create new GraphBLAS objects
using one of two methods:  (1) per-object with methods such as
\verb'GxB_Matrix_new_arena' (see the the next section) or (2) for all objects
using the global context or a specific \verb'Context' object for any given
thread in the user application (described below). 

For method (2), the get/set options are used.  To set the global default
arenas, use:

    {\footnotesize
    \begin{verbatim}
    GrB_set (GrB_GLOBAL, arena, GxB_ARENA_DATA) ;
    GrB_set (GrB_GLOBAL, arena, GxB_ARENA_HEADER) ; \end{verbatim}}

To set the arenas via a \verb'Context' object, use:

    {\footnotesize
    \begin{verbatim}
    GxB_set (Context, arena, GxB_ARENA_DATA) ;
    GxB_set (Context, arena, GxB_ARENA_HEADER) ; \end{verbatim}}

Once either of the above settings are used, all subsequent newly-created
objects will be in the new arenas.  For the \verb'Context' object, the
\verb'Context' must also be engaged.  See
Sections~\ref{context}, \ref{get_set_context}, and \ref{get_set_global}
for details.

These settings can be queried via \verb'GrB_get':

    {\footnotesize
    \begin{verbatim}
    int32_t arena ;
    GrB_get (GrB_GLOBAL, &arena, GxB_ARENA_DATA) ;
    GrB_get (GrB_GLOBAL, &arena, GxB_ARENA_HEADER) ;
    GxB_get (Context, &arena, GxB_ARENA_DATA) ;
    GxB_get (Context, &arena, GxB_ARENA_HEADER) ; \end{verbatim}}

This is the simplest method to use different arenas, but it does not allow
control on a per-object basis.

%-------------------------------------------------------------------------------
\subsection{Using arenas via per-object methods}
%-------------------------------------------------------------------------------

{\footnotesize
\begin{verbatim}

//      GxB_Descriptor_new_arena
//      GxB_Type_new_arena
//      GxB_UnaryOp_new_arena
//      GxB_BinaryOp_new_arena
//      GxB_IndexBinaryOp_new_arena
//      GxB_IndexUnaryOp_new_arena
//      GxB_Monoid_new_arena*
//      GxB_Monoid_terminal_new_arena*
//      GxB_Semiring_new_arena
//      GxB_Context_new_arena
//      GxB_Container_new_arena
//      GxB_Iterator_new_arena

//      GxB_Scalar_new_arena
//      GxB_Vector_new_arena
//      GxB_Matrix_new_arena

//      GxB_Scalar_dup_arena
//      GxB_Vector_dup_arena
//      GxB_Matrix_dup_arena
//      GxB_Matrix_split_arena
//      GxB_Matrix_diag_arena
//      GxB_Matrix_serialize_arena
//      GxB_Vector_serialize_arena
//      GxB_Matrix_deserialize_arena
//      GxB_Vector_deserialize_arena
//      GxB_Matrix_reshapeDup_arena

    to change both header & data arena: (not lazy; always moves data if needed)
    GxB_Matrix_set_arenas (&A, header_arena, data_arena) ;
    GxB_Vector_set_arenas (&V, header_arena, data_arena) ;
    GxB_Scalar_set_arenas (&S, header_arena, data_arena) ;

\end{verbatim}}

%-------------------------------------------------------------------------------

{\footnotesize
\begin{verbatim}
FIXME arena: add this to user guide:

    int arena ;

    changes the data arena only: (not lazy; always moves data if needed)
    but could be done as a lazy, leaving it as pending work:
    GrB_Matrix_set_INT32 (A, arena, GxB_ARENA_DATA) ;
    GrB_Vector_set_INT32 (V, arena, GxB_ARENA_DATA) ;
    GrB_Scalar_set_INT32 (S, arena, GxB_ARENA_DATA) ;

    get the data or header arena of a matrix:
    GrB_Matrix_get_INT32 (A, &arena, GxB_ARENA_DATA) ;
    GrB_Vector_get_INT32 (V, &arena, GxB_ARENA_DATA) ;
    GrB_Scalar_get_INT32 (S, &arena, GxB_ARENA_DATA) ;
    GrB_Matrix_get_INT32 (A, &arena, GxB_ARENA_HEADER) ;
    GrB_Vector_get_INT32 (V, &arena, GxB_ARENA_HEADER) ;
    GrB_Scalar_get_INT32 (S, &arena, GxB_ARENA_HEADER) ;

    get the arena for all other objects:
    GrB_*_get_INT32 (object, &arena, GxB_ARENA_HEADER) ;

\end{verbatim}}

